\documentclass[11pt]{article}
\usepackage{fontspec}
\usepackage[french,english]{babel}
\usepackage{amsmath}
\usepackage{amssymb}
\usepackage{amsthm}
\usepackage{mathrsfs}
\usepackage{tikz-cd}
\usepackage{enumitem}

\frenchsetup{AutoSpacePunctuation=false}
\setlength{\parindent}{0pt}
\setlength{\parskip}{0.5\baselineskip}

\newcommand{\printedpage}[1]{\marginpar{\footnotesize p.~#1}}
\newcommand{\msp}[1]{\marginpar{\footnotesize [#1]}}
\newcommand{\und}[1]{\underline{\mathrm{#1}}}
\newcommand{\HOT}{\mathrm{HOT}}
\newcommand{\hot}{\mathrm{hot}}
\newcommand{\Cat}{\mathrm{Cat}}

\begin{document}

\printedpage{74}
% head: A. Grothendieck : Dérivateurs, Ch. VI, Édition des Universités de Montpellier II et Paris VII
% foot: Version 9 novembre 2001

\msp{page 119}
A similar remark applies to $\und{Fibhot}\,S$, where $W_S^{\mathrm{d}} = W_S^{\mathrm{f}}$ as well. The base change functor induces
\[
\und{Fibhot}\,S \longrightarrow \und{Fibhot}\,S', \qquad X \longmapsto X \times_S S'/S',
\]
compatible with the relevant relative localizers (all those considered so far coincide on these categories, under the hypotheses made on $W$), whence
\[
f^* : (\und{Fibhot}\,S)W_S^{-1} \longrightarrow (\und{Fibhot}\,S')W_{S'}^{-1}.
\]
The localized categories this time are not the $\HOT_W$, but the $\HOT_W^{\mathrm{lc}}$, so we find a nice description of
\[
f^* : \HOT_W^{\mathrm{lc}}(S) \longrightarrow \HOT_W^{\mathrm{lc}}(S'),
\]
still more advantageous than in terms of $\und{Liss}_0S$, $\und{Liss}_0S'$, since applying to somewhat larger categories (NB $\und{Liss}_0 \subseteq \und{Fibhot}$). Let us therefore note the

\textbf{Corollary 1.} \textit{The formula of proposition 1 (p. 118) remains valid if $X$ is assumed} Hot-\textit{fibrant over $S$ (instead of smooth over $S$).} (One then finds relative homotopy types \textit{locally constant} over $S$, $S'$.)

We thus have a description of $f^*$ which applies to all the categories on the left-hand side
\msp{page 120}
of the diagram \texttt{[page]} 76, with the sole exception of $\Cat/S$ itself. If $X$ is in $\Cat/S$, without being either smooth or Hot-fibered over $S$, how to find the inverse image in $\HOT_W(S')$? We have the choice between $\Phi$ and $C$ to improve the situation! We know that $X$ defines the same class in $\HOT_W(S)$ as $\Phi(X)$, $C(X)$, which are, the one smooth (and even split fibered), the other Hot-fibered (and even perfect) over $S$. Thus

\textbf{Corollary 2.} \textit{If $X$ in} $\Cat/S$ \textit{is arbitrary, then}
\[
f^*(\hot_{W,S}(X)) \simeq \hot_{W,S'}(\Phi(X) \times_S S') \simeq \hot_{W,S'}(C(X) \times_S S').
\]

\bigskip

One may expect, on the other hand, that for certain particular base changes, the formula of proposition 1 is valid without restriction on $X$ in $\Cat/S$. It must come to $\pm$ the same thing that the inverse image functor
\[
X \longmapsto X \times_S S'
\]
be compatible with the localizers $W_S^{\mathrm{d}}$, $W_{S'}^{\mathrm{d}}$. The condition of $W$-equivalence ``colocal over $S$\,'' must be viewed as a condition of a ``homological'' nature, for the derivator HOT. It is associated, for an arbitrary derivator $\mathbf{D}$, with the condition (for $f : X \longrightarrow Y$ in $\Cat/S$) that the $s\backslash X \longrightarrow s\backslash Y$ are $\mathbf{D}$-equivalences,
\printedpage{75}
% head: A. Grothendieck : Dérivateurs, Ch. VI, Édition des Universités de Montpellier II et Paris VII
% foot: Version 9 novembre 2001
\msp{page 121}
which is written here elegantly in terms of relative homology, by the condition that $\forall\ a \in \mathcal{A} = \mathbf{D}(e)$, the arrow induced by $f$ % [sic] induit -> induite, accord avec « la flèche »
\[
\underbrace{H_\bullet^{\mathbf{D}}(X/S, a_X)}_{\overset{\text{def}}{=}(p_X)_!(a_X)} \longrightarrow \underbrace{H_\bullet^{\mathbf{D}}(Y/S, a_Y)}_{\overset{\text{def}}{=}(p_Y)_!(a_Y)}
\]
be \texttt{[an]} isomorphism. (Use the ``calculation'' of the relative homology sheaves $(p_X)_!(a_X)$, $(p_Y)_!(a_Y)$, by the characterization of the fibers in Der 4 bis). Now, the $\mathbf{D}$-\textit{proper} base changes $S' \longrightarrow S$ commute with the $p_!$, so that we find the

\textbf{Lemma.} \textit{The notion of} $\mathbf{D}$\textit{-colocal equivalence over $S$ in} $\Cat/S$ \textit{is stable under every base change $S' \longrightarrow S$ that is} $\mathbf{D}$\textit{-proper. (Whatever the derivator} $\mathbf{D}$\textit{.)}

It will be recalled that the $s\backslash S \longrightarrow S$ are precisely proper base changes! One % [sic] rappelera -> rappellera, orthographe; des base -> de base, accord
can say that $\mathbf{D}$-colocal equivalence can be described by requiring that $f : X \longrightarrow Y$ be transformed into $W$-equivalences by certain types of base change $S' \longrightarrow S$, choosing at will one of the three types

\begin{itemize}
\item[a)] The base changes $s\backslash S \longrightarrow S$ (which gives back the definition we had adopted).
\item[b)] The base changes that are local co-isomorphisms.
\item[c)] Proper base changes.
\item[d)] $\mathbf{D}$-proper base changes.
\end{itemize}

We can apply these results only to the derivator
\msp{page 122}
HOT, for lack of having yet constructed it, and verified, the relevant axioms. We can simply state % [sic] verifié -> vérifié, accent
the

\textbf{Conjecture} (sure!) : \textit{Let $p : S' \longrightarrow S$ be a proper arrow (understood:} HOT\textit{-proper) in} Cat. \textit{Then the base change functor}
\[
\Cat/S \longrightarrow \Cat/S', \qquad X \longmapsto X \times_S S'/S'
\]
\textit{is compatible with the localizers $W_S^{\mathrm{d}}$, $W_{S'}^{\mathrm{d}}$. (For a smooth functor $S' \longrightarrow S$, base change would be compatible with $W_S^{\mathrm{g}}$, $W_{S'}^{\mathrm{g}}$.)} (\footnote[61]{NB This ``conjecture'' is proved without a lemma on page 168-169, without needing anything that is not % [sic] prouvé -> prouvée, accord avec « conjecture »
already known here. Cf. theorem 4, case b) p. 162.}).

But for every base change $S' \longrightarrow S$ (proper or not) which has the preceding property, % [sic] pout -> pour, faute de frappe
hence which induces a functor
\[
\underbrace{(\Cat/S)(W_S^{\mathrm{d}})^{-1}}_{\approx\,\HOT_W(S)} \longrightarrow \underbrace{(\Cat/S')(W_{S'}^{\mathrm{d}})^{-1}}_{\approx\,\HOT_W(S')}\ ,
\]
\printedpage{76}
% head: A. Grothendieck : Dérivateurs, Ch. VI, Édition des Universités de Montpellier II et Paris VII
% foot: Version 9 novembre 2001
it is immediate, using the comparison theorem (e.g. with $\und{Liss}\,S$, $\und{Liss}\,S'$) over $S$ and over $S'$, that this functor is none other (up to canonical isomorphism) than the inverse image functor $f^*$ for relative homotopy types. So let us state the % [sic] énoncons -> énonçons, cédille

\textbf{Corollary 3.} \textit{The formula of proposition 1 remains valid without any hypothesis whatever on $X$ in} $\Cat/S$\textit{, provided that the base change arrow $f : S' \longrightarrow S$ is such that the functor $X \longmapsto X \times_S S'$ from} $\Cat/S$ \textit{to} $\Cat/S'$
\msp{page 123}
\textit{takes colocal $W$-equivalences over $S$ to colocal $W$-equivalences over $S'$.} (Which should be the case if (and only if?) $f$ is $W$-proper, hence also if $f$ is (``absolutely'') proper (\footnote[62]{cf. I, pages 55-59.}).)

So far we have worked only with the localizers $W_S^{\mathrm{d}}$, $W_{S'}^{\mathrm{d}}$, to the exclusion of $W_S$, $W_{S'}$. (Except on the categories contained in $\und{Fibhot}\,S$, where the two coincide.) When one makes no draconian hypothesis on $p : S' \longrightarrow S$, of the Hot-fibration type, one has no chance of $X \longmapsto X \times_S S'$ taking coarse $W$-equivalences into likewise, hence no chance of catching the inverse image of the ``coarse'' relative homotopy type of $X$, by taking that of $X' = X \times_S S'$. This is tied to the fact that for a general $S' \longrightarrow S$, the localization functors
\[
\HOT_W(S) \xrightarrow{\rho_S} \HOT_W^{\mathrm{lc}}(S), \qquad \HOT_W(S') \xrightarrow{\rho_{S'}} \HOT_W^{\mathrm{lc}}(S')
\]
have no desire to commute with the inverse image functor $f^*$ on the $\HOT_W$. We may hazard the

\textbf{Proposition 2.} \textit{Equivalent conditions on the functor $S' \longrightarrow S$ :}

\begin{itemize}
\item[\textit{a)}] \textit{The diagram of functors}
\[
\begin{tikzcd}
\HOT_W(S) \arrow[r, "\rho_S"] \arrow[d, "f^*"] & \HOT_W^{\mathrm{lc}}(S) \arrow[d, "f^*_{\mathrm{lc}}"] \\
\HOT_W(S') \arrow[r, "\rho_{S'}"] & \HOT_W^{\mathrm{lc}}(S')
\end{tikzcd}
\]
\textit{commutes (up to isomorphism of composite functors).}

\msp{page 124}

\item[\textit{b)}] \textit{The base change functor} % [sic] changment -> changement, faute de frappe
\[
\und{Liss}\,S \longrightarrow \und{Liss}\,S', \qquad X \longmapsto X \times_S S'/S'
\]
\textit{takes $W_S$ into $W_{S'}$.}
\printedpage{77}
% head: A. Grothendieck : Dérivateurs, Ch. VI, Édition des Universités de Montpellier II et Paris VII
% foot: Version 9 novembre 2001
\item[\textit{b$'$)}] \textit{Analogous statement for} $\und{Fib}\,S \longrightarrow \und{Fib}\,S'$\textit{.}
\item[\textit{b$''$)}] \textit{Analogous statement for} $\und{Fibsc}\,S \longrightarrow \und{Fibsc}\,S'$\textit{.}
\end{itemize}

(\footnote[63]{Not proved, only that b, b$'$, b$''$ imply a). For the reverse implication, one would need to know that every arrow of $\mathcal{F}'$ which becomes \texttt{[an]} isomorphism in $\mathcal{F}'W_{S'}^{-1}$, is in $W_{S'}$ (strong saturation of $W_{S'}$ in $\mathcal{F}'$), something I do not know. But I believe this must be true.}, \footnote[64]{\textbf{NB} If we assume $p$ proper, we should also be able to take $\mathcal{F} = \Cat/S$.}).

\textsc{Proof.} If $\mathcal{F}$ is one of the three categories $\und{Liss}\,S$, $\und{Fib}\,S$ or $\und{Fibsc}\,S$, the functor $\rho_S$ can be described as the canonical localization functor
\[
\mathcal{F}(W_S^{\mathrm{d}})^{-1} \longrightarrow \mathcal{F}W_S^{-1}
\]
(that would work for any full subcategory $\mathcal{F}$ of $\Cat/S$ which contains $\und{Fib}\,S$). Now if b), b$'$) or b$''$) is satisfied, we find a commutative diagram of functors induced by base change
\[
\begin{tikzcd}
\mathcal{F}(W_S^{\mathrm{d}})^{-1} \arrow[r, "\rho_S"] \arrow[d, "f^*_{W_S^{\mathrm{d}}}"'] & \mathcal{F}(W_S)^{-1} \arrow[d, "f^*_{W_S}"] \\
\mathcal{F}'(W_{S'}^{\mathrm{d}})^{-1} \arrow[r, "\rho_{S'}"] & \mathcal{F}'(W_{S'})^{-1},
\end{tikzcd}
\]
with the proviso, it is true, in the case $\mathcal{F} = \Cat/S$, that the inverse image functor is also compatible with $W_S^{\mathrm{d}}$, $W_{S'}^{\mathrm{d}}$ (which is automatic if $\mathcal{F} = \und{Fib}\,S$ or $\und{Fibsc}\,S$). We know that $f^*_{W_S^{\mathrm{d}}}$ describes the functor $f^*$ on relative homotopy types, and to prove a) as a consequence
\msp{page 125}
of b) (strengthened), b$'$ or b$''$, it remains to make sure that the functor $f^*_{W_S}$ likewise describes the functor induced by the preceding one on locally constant relative homotopy types. Now we have the canonical inclusions $\mathcal{F}_0 \subseteq \mathcal{F}$, $\mathcal{F}_0' \subseteq \mathcal{F}'$, and on $\mathcal{F}_0$, $\mathcal{F}_0'$ \texttt{[we have]} $W_S = W_S^{\mathrm{d}} = W_S^{\mathrm{f}}$, and the categories are OK for describing $\HOT_W^{\mathrm{lc}}(S)$ and $\HOT_W^{\mathrm{lc}}(S')$, and the functor $f^*$ induced between them by inverse images (\footnote[65]{NB $(\Cat/S)_0 = \und{Fibhot}\,S$, for the other cases, it is clear.}) (\textbf{NB} The base change functor $X \longmapsto X \times_S S'$ is compatible with the localizers on these categories, for it is trivial in the form $W_S^{\mathrm{f}}$.) We then have, by inverse images $X \longmapsto X \times_S S'$, a commutative diagram
\[
\begin{tikzcd}
\mathcal{F}_0(W_S)^{-1} \arrow[r, "\approx"] \arrow[d, "(f^*_{W_S})_0"'] & \mathcal{F}(W_S)^{-1} \arrow[d, "f^*_{W_S}"] \\
\mathcal{F}_0'(W_{S'})^{-1} \arrow[r, "\approx"] & \mathcal{F}'(W_{S'})^{-1},
\end{tikzcd}
\]
\printedpage{78}
% head: A. Grothendieck : Dérivateurs, Ch. VI, Édition des Universités de Montpellier II et Paris VII
% foot: Version 9 novembre 2001
where the horizontal arrows are \textit{equivalences} of categories, by theorem-scholie 1 \textcircled{I}~. This proves that $f^*_{W_S}$ too is good for calculating the inverse images of relative homotopy types.

We have proved that b, b$'$, b$''$ each imply a), let us prove that a) implies them. Now a) says exactly that the inverse images $f^*$ of the elements of $\HOT_W^{\mathrm{lc}}(S)$
\msp{page 126}
are, up to isomorphism, obtained by the inverse image functor on the whole of $\mathcal{F}$, interpreting $\HOT_W^{\mathrm{lc}}(S)$ as $\mathcal{F}W_S^{-1}$. But this does indeed imply that $\mathcal{F} \longrightarrow \mathcal{F}'$ must be compatible with the localizers. But I'm getting carried away - one only finds that if $u : X \longrightarrow Y$ in $\mathcal{F}$ is in $W_S$, then the arrow $u' : X' \longrightarrow Y'$ gives an isomorphism in $\mathcal{F}'(W_{S'})^{-1}$. I do not know, on the other hand, whether this implies that $u'$ is in $W_{S'}$.

\textbf{Corollary 1.} \textit{Suppose that $p : S' \longrightarrow S$ is a} Hot\textit{-fibration. Then the inverse image functor $X \longmapsto X \times_S S'$ from} $\Cat/S$ \textit{to $\mathit{Cat}/S'$ maps $W_S$ into $W_{S'}$, hence the functor $p^* : \HOT_W(S) \longrightarrow \HOT_W(S')$ commutes with the localization functors $\rho_S$, $\rho_{S'}$ in these categories.} (\textbf{NB} One must make the appropriate hypothesis on $W$, namely L6 bis page 44. It is OK for $W_\infty$.)

The fact that $X \longmapsto X \times_S S'$ is compatible with $W_S$, $W_{S'}$ is a tautology, in view of hypothesis L6 bis. It does not seem that it is also compatible with $W_S^{\mathrm{d}}$, $W_{S'}^{\mathrm{d}}$ (I should make a counterexample, with of course $S'$ not proper over $S$). So we are under the conditions of b), which implies a).

\msp{page 127}

I do not feel capable of exhibiting a functor other than a Hot-fibration (or at least a $W$-fibration), for which the inverse image functor on $\HOT_W$ is compatible with the localization functors $\rho_S$, $\rho_{S'}$. Of course, if $X$ is arbitrary in $\Cat/S$, nothing prevents us from looking at its localized object for $W_S$, i.e. $\rho_S(\hot_{W,S}(X))$, and asking how to find its inverse image. If $p : S' \longrightarrow S$ is a Hot-fibration, then it suffices to take $\hot_{W,S'}(X')$, $X' = S \times_S S'$. In any other case (I mean, without a hypothesis on $X/S$, more % [sic] X' = S \times_S S' -> X' = X \times_S S', cf. la définition de X' plus haut
precisely without assuming that $X$ is itself Hot-fibered over $S$), this formula is hopelessly botched. We must, to calculate $f^*\rho_S(\hot_{W,S}(X))$, ``calculate'' $\rho_S(\hot_{W,S}(X))$ by ``going up to the upper floor'', as $\hot_{W,S}(C(X))$. We then find the inverse image $\hot_{W,S'}(C(X) \times_S S')$ (by \texttt{[the]} corollary 1, p. 119).

Now that the contravariant law of $\HOT_W(S)$ in $S \in \mathrm{Ob}\,\Cat$ seems well understood, it is a matter of seeing that we do indeed have a derivator HOT. I do not expect the $\HOT_W^{\mathrm{lc}}(S)$
\msp{page 128}
to define a derivator as well, but I would have to make sure of this expressly. (It is for the existence of $f_!$, $f_*$ that I must look.) % [sic] je doit -> je dois, conjugaison

To tell the truth, one would have to, to do things properly, look also at the 2-\textit{functor} structure of
\[
S \longmapsto \HOT_W(S) \qquad \Cat^o \longrightarrow \Cat,
\]

\end{document}
